%===========================================================================


\section{Model Diagnostics}
\label{sec:model_diagnostics}
%===========================================================================

Two checks on the comparator models are collected here: a test of the Cox
comparator's proportional-hazards assumption, and the predictors the
gradient-boosted trees used, which bear on Sections~\ref{sec:ablation}
and~\ref{sec:h2}.

\subsection{Proportional-Hazards Testing \texorpdfstring{[C]}{[C]}}
\label{sec:schoenfeld}

The Schoenfeld residual test on the Cox comparator (Variant B) rejects
proportional hazards ($\chi^2 = \pcSchoenfeldChisq$, df = \pcSchoenfeldDf,
$p = \pcSchoenfeldP$). Because the comparator uses hourly covariates but
time-constant coefficients, the test rejects the hazard-ratio structure
(Section~\ref{sec:cox_spec}).
\pcSchoenfeldNViolate{} of the \pcSchoenfeldNCovariates{} covariates violate
the assumption individually at $\alpha = 0.05$, led by GCS, bilirubin,
respiratory rate and creatinine. These are per-covariate tests reported without
multiplicity adjustment; the global test alone is sufficient for the
pre-registered check, and the smallest p-values survive Bonferroni correction
over the covariate set by many orders of magnitude. The Cox coefficients
therefore cannot be read as time-constant hazard ratios on these data.

\subsection{GBT Feature Importance \texorpdfstring{[E]}{[E]}}
\label{sec:feature_importance}

\begin{table}[H]
    \centering
    \begin{tabular}{lrrr}
        \toprule
        \textbf{Feature}         & \textbf{Gain} & \textbf{Cover} & \textbf{Frequency} \\
        \midrule
        Temperature ($^\circ$C)  & \pcGbtGainTempC & \pcGbtCoverTempC & \pcGbtFreqTempC \\
        Resp.\ rate slope (6\,h) & \pcGbtGainRespRateSlope & \pcGbtCoverRespRateSlope & \pcGbtFreqRespRateSlope \\
        MAP slope (6\,h)         & \pcGbtGainMapSlope & \pcGbtCoverMapSlope & \pcGbtFreqMapSlope \\
        GCS                      & \pcGbtGainGcs & \pcGbtCoverGcs & \pcGbtFreqGcs \\
        Heart-rate slope (6\,h)  & \pcGbtGainHrSlope & \pcGbtCoverHrSlope & \pcGbtFreqHrSlope \\
        Creatinine               & \pcGbtGainCreatinine & \pcGbtCoverCreatinine & \pcGbtFreqCreatinine \\
        WBC                      & \pcGbtGainWbc & \pcGbtCoverWbc & \pcGbtFreqWbc \\
        Heart rate (level)       & \pcGbtGainHr & \pcGbtCoverHr & \pcGbtFreqHr \\
        Lactate                  & \pcGbtGainLactate & \pcGbtCoverLactate & \pcGbtFreqLactate \\
        MAP (level)              & \pcGbtGainMap & \pcGbtCoverMap & \pcGbtFreqMap \\
        \bottomrule
    \end{tabular}
    \caption{Top-10 GBT features by information gain (Variant B). Three of the
        five most informative predictors are six-hour trend features.}
    \label{tab:feature_importance}
\end{table}

The prominence of six-hour slope features is consistent with Lauritsen et al.'s
\cite{lauritsen2021framing} concern: the GBT model draws much of its signal from
trends, so its performance depends on the length of the history window used for
feature construction, a choice rarely reported or varied in the comparable
literature. Heart rate contributes through both its level and its slope.
